\begin{definition}[Графовое число Рамсея]
Для графов \(G_1,G_2\) число \(R(G_1,G_2)\) — наименьшее \(t\) такое, что при любой
раскраске рёбер \(K_t\) в два цвета найдётся подграф цвета 1, изоморфный \(G_1\),
или подграф цвета 2, изоморфный \(G_2\). Подграф \emph{не обязан быть индуцированным}:
нужны лишь рёбра \(G_i\) нужного цвета, остальные рёбра на тех же вершинах любые.
Аналогично для \(k\) цветов: \(R(G_1,\ldots,G_k)\).
\end{definition}
\begin{theorem}[Существование]
\[
R(G_1,\ldots,G_k)\le R\bigl(|G_1|,\ldots,|G_k|\bigr),
\]
где \(|G_i|=v(G_i)\). В частности, графовые числа Рамсея конечны.
\end{theorem}
\begin{proof}
Пусть \(t=R(|G_1|,\ldots,|G_k|)\) и дана раскраска \(K_t\). По определению найдётся \(i\) и
клика \(K_{|G_i|}\) цвета \(i\). Граф \(G_i\) на \(|G_i|\) вершинах — подграф полного графа
на стольких же вершинах, поэтому в этой клике есть копия \(G_i\), все рёбра которой
цвета \(i\). Конечность следует из билета~\ref{ticket:10}.
\end{proof}
\begin{cor}
\(R(K_{n_1},\ldots,K_{n_k})=R(n_1,\ldots,n_k)\): для клик неравенство превращается в
определение.
\end{cor}
\begin{graphillustration}{Слева: одноцветная клика $K_4$ содержит любой граф на 4 вершинах, например цикл $C_4$ (толстые рёбра); поэтому $R(C_4,\cdot)\le R(4,\cdot)$. Справа: в любой 2-раскраске $K_3$ два ребра одного цвета образуют $P_3$.}{fig:11-embed}
\begin{tikzpicture}[font=\small]
\begin{scope}
\node[gi vertex] (a) at (0,0) {};
\node[gi vertex] (b) at (1.6,0) {};
\node[gi vertex] (c) at (1.6,1.6) {};
\node[gi vertex] (d) at (0,1.6) {};
\draw[gi c1] (a) -- (c); \draw[gi c1] (b) -- (d);
\draw[gi match] (a) -- (b) -- (c) -- (d) -- (a);
\node[gi note] at (0.8,-0.6) {$C_4\subseteq K_4$ цвета 1};
\end{scope}
\begin{scope}[shift={(4.2,0)}]
\node[gi vertex] (x) at (0,0) {};
\node[gi vertex] (y) at (1.8,0) {};
\node[gi vertex] (z) at (0.9,1.5) {};
\draw[gi c1,line width=2pt] (x) -- (z); \draw[gi c1,line width=2pt] (z) -- (y);
\draw[gi c2] (x) -- (y);
\node[gi note] at (0.9,-0.6) {одноцветный $P_3$};
\end{scope}
\end{tikzpicture}
\end{graphillustration}

\begin{example}[Графовые числа бывают много меньше]
Несколько простых значений:
\begin{itemize}\tightlist
\item \(R(P_3,P_3)=3\) (\(P_3\) — путь из двух рёбер), тогда как \(R(3,3)=6\).
  В \(K_3\) из трёх рёбер два одного цвета, а любые два ребра треугольника
  имеют общую вершину — одноцветный \(P_3\). В \(K_2\) одно ребро, \(P_3\) нет.
\item \(R(K_2,G)=v(G)\) при \(v(G)\ge2\): раскраска без ребра цвета 1 — сплошь цвета 2, и
  \(G\subseteq K_t\iff t\ge v(G)\).
\item \(R(T_n,K_m)=(m-1)(n-1)+1\) для любого дерева \(T_n\) — билет~\ref{ticket:12}.
\end{itemize}
\end{example}
\begin{tcolorbox}[colback=cbOrangeBack,colframe=cbOrange!45!cbGray,
  borderline west={1.6pt}{0pt}{cbOrange}]
\textbf{Индуцированный случай (только формулировка).} Граф \(G\) называется
\(H\)-рамсеевским, если при любой раскраске рёбер \(G\) в два цвета есть одноцветный
\emph{индуцированный} подграф \(G\), изоморфный \(H\). Для любого \(H\) такой \(G\)
существует (индуцированная теорема Рамсея: Деуб, Эрдёш–Хайнал–Поша, Рёдль, 1970-е),
но доказательство значительно сложнее и в программу не входит. Для клик
«индуцированный» ничего не меняет — это обычная теорема Рамсея.
\end{tcolorbox}
